\documentclass[12pt]{article}
\usepackage[utf8]{inputenc}
\usepackage[T1]{fontenc}
\usepackage{hyperref}
\title{Tichurong Language Resources for K-12 Curriculum Development}
\author{Compiled by Claude Sonnet 5 (claude-sonnet-5) for Andr\'as Kornai}
\date{}
\begin{document}
\maketitle
\section{Introduction}
Tichurong (ISO 639-3: tcn; Glottocode tich1238) is spoken in the Tichurong valley, Dolpa district, Karnali Province, Nepal. Glottolog classifies it as Tibetic, in a Dolpo-Tichurong group. Speaker estimates are about 2,400 (older Ethnologue-based figure) and about 2,700 across eighteen villages (a more recent source). Other names include Poike, Poinke, Rongke and Tichurongke. It is unwritten. It is used in everyday life, though usage varies with age and gender, and most residents also speak Nepali and Tibetan. The valley's other native language, Kaike, is not Tibetic.
\section{Primary Educational Resources}
I found no primers or school materials. The language is unwritten and not known to be taught in schools.
\section{Linguistic and Reference Materials}
\begin{itemize}
\item Language Documentation and Description journal, ``Tichurong (Nepal) -- Language Snapshot''. \url{https://lddjournal.org/article/id/1257}
\item Same snapshot, PDF download. \url{https://www.lddjournal.org/article/1257/galley/2500/download}
\item Honda, Isao (2018), preliminary report on Tichurong, ICSTLL 51 proceedings (Kyoto University repository). \url{https://repository.kulib.kyoto-u.ac.jp/dspace/bitstream/2433/235280/1/proc_icstll51_25.pdf}
\item Glottolog reference record, Honda 2018. \url{https://glottolog.org/resource/reference/id/572699}
\item Glottolog reference record, Budha, Daurio and Turin 2021. \url{https://glottolog.org/resource/reference/id/602676}
\item Numerals page from the Max Planck EVA numerals project. \url{https://lingweb.eva.mpg.de/channumerals/Tichurong.htm}
\item Wikipedia, Tichurong language. \url{https://en.wikipedia.org/wiki/Tichurong_language}
\item Glottolog languoid page. \url{https://glottolog.org/resource/languoid/id/tich1238.bigmap.html}
\end{itemize}
\section{Religious and Community Resources}
I found none in the search results.
\section{Digital and Audio Resources}
\begin{itemize}
\item Wikitongues page for the language. \url{https://wikitongues.org/languages/dre/}
\end{itemize}
I did not verify what this page contains. The URL code (dre) may point to a different language, so check it before use.
\section{Curriculum Development Context}
The language is stable and used by the whole community, but there is no standard orthography and no known school use. The documentation found is a language snapshot, a preliminary survey report and a 2021 reference by Budha, Daurio and Turin. A 2018 sociolinguistic survey through Nepal's Language Commission was mentioned in the search summary, and the language was reported as one of 15 flagged for further research by the Central Department of Linguistics. A mother-tongue curriculum would first need an agreed writing system and teaching materials, developed with the community and Nepali education authorities.
\end{document}